\begin{figure*}[t]
\centering
\begin{tikzpicture}
\draw[ge] (4.917,2.297) -- (4.335,0.886);
\draw[ge] (4.917,2.297) -- (3.738,2.659);
\draw[ge] (4.917,2.297) -- (4.422,3.350);
\draw[ge] (4.917,2.297) -- (4.054,1.780);
\draw[ge] (2.718,2.836) -- (3.003,1.824);
\draw[ge] (2.718,2.836) -- (3.738,2.659);
\draw[ge] (2.718,2.836) -- (4.422,3.350);
\draw[ge] (2.718,2.836) -- (4.054,1.780);
\draw[ge] (2.718,2.836) -- (1.997,2.076);
\draw[ge] (3.003,1.824) -- (4.335,0.886);
\draw[ge] (3.003,1.824) -- (3.738,2.659);
\draw[ge] (3.003,1.824) -- (4.054,1.780);
\draw[ge] (3.003,1.824) -- (3.285,0.894);
\draw[ge] (4.335,0.886) -- (4.054,1.780);
\draw[ge] (3.738,2.659) -- (4.422,3.350);
\draw[ge] (3.738,2.659) -- (4.054,1.780);
\draw[ge] (4.422,3.350) -- (4.054,1.780);
\draw[ge] (4.054,1.780) -- (3.285,0.894);
\draw[ge] (3.285,0.894) -- (2.049,1.109);
\draw[ge] (2.049,1.109) -- (1.997,2.076);
\draw[ge] (2.049,1.109) -- (0.937,1.503);
\draw[ge] (1.997,2.076) -- (1.403,2.888);
\node[gn, fill=pc0, text=white] (P1n1) at (4.917,2.297) {1};
\node[gn, fill=pc0, text=white] (P1n2) at (2.718,2.836) {2};
\node[gn, fill=pc0, text=white] (P1n3) at (3.003,1.824) {3};
\node[gn, fill=pc0, text=white] (P1n4) at (4.335,0.886) {4};
\node[gn, fill=pc0, text=white] (P1n5) at (3.738,2.659) {5};
\node[gn, fill=pc0, text=white] (P1n6) at (4.422,3.350) {6};
\node[gn, fill=pc0, text=white] (P1n7) at (4.054,1.780) {7};
\node[gn, fill=pc0, text=white] (P1n8) at (3.285,0.894) {8};
\node[gn, fill=pc1, text=white] (P1n9) at (2.049,1.109) {9};
\node[gn, fill=pc1, text=white] (P1n10) at (1.997,2.076) {10};
\node[gn, fill=pc1, text=white] (P1n11) at (1.403,2.888) {11};
\node[gn, fill=pc2, text=black] (P1n12) at (0.937,1.503) {12};
\node[ptitle] at (2.583,3.600) {(a) Stage 1 result};
\node[pstat] at (2.583,0.000) {3 comm., $Q{=}0.168$, NMI 0.872};
\node[barrow, rotate=0, minimum height=9mm] at (5.817,1.800) {};
\node[note, anchor=south] at (5.817,2.080) {Stage 2};
\draw[ge] (11.383,2.297) -- (10.802,0.886);
\draw[ge] (11.383,2.297) -- (10.204,2.659);
\draw[ge] (11.383,2.297) -- (10.889,3.350);
\draw[ge] (11.383,2.297) -- (10.521,1.780);
\draw[ge] (9.185,2.836) -- (9.470,1.824);
\draw[ge] (9.185,2.836) -- (10.204,2.659);
\draw[ge] (9.185,2.836) -- (10.889,3.350);
\draw[ge] (9.185,2.836) -- (10.521,1.780);
\draw[ge] (9.185,2.836) -- (8.464,2.076);
\draw[ge] (9.470,1.824) -- (10.802,0.886);
\draw[ge] (9.470,1.824) -- (10.204,2.659);
\draw[ge] (9.470,1.824) -- (10.521,1.780);
\draw[ge] (9.470,1.824) -- (9.751,0.894);
\draw[ge] (10.802,0.886) -- (10.521,1.780);
\draw[ge] (10.204,2.659) -- (10.889,3.350);
\draw[ge] (10.204,2.659) -- (10.521,1.780);
\draw[ge] (10.889,3.350) -- (10.521,1.780);
\draw[ge] (10.521,1.780) -- (9.751,0.894);
\draw[ge] (9.751,0.894) -- (8.515,1.109);
\draw[ge] (8.515,1.109) -- (8.464,2.076);
\draw[ge] (8.515,1.109) -- (7.404,1.503);
\draw[ge] (8.464,2.076) -- (7.869,2.888);
\node[ring] at (9.751,0.894) {};
\node[gn, fill=pc0, text=white] (P2n1) at (11.383,2.297) {1};
\node[gn, fill=pc0, text=white] (P2n2) at (9.185,2.836) {2};
\node[gn, fill=pc0, text=white] (P2n3) at (9.470,1.824) {3};
\node[gn, fill=pc0, text=white] (P2n4) at (10.802,0.886) {4};
\node[gn, fill=pc0, text=white] (P2n5) at (10.204,2.659) {5};
\node[gn, fill=pc0, text=white] (P2n6) at (10.889,3.350) {6};
\node[gn, fill=pc0, text=white] (P2n7) at (10.521,1.780) {7};
\node[gn, fill=pc1, text=white] (P2n8) at (9.751,0.894) {8};
\node[gn, fill=pc1, text=white] (P2n9) at (8.515,1.109) {9};
\node[gn, fill=pc1, text=white] (P2n10) at (8.464,2.076) {10};
\node[gn, fill=pc1, text=white] (P2n11) at (7.869,2.888) {11};
\node[gn, fill=pc1, text=white] (P2n12) at (7.404,1.503) {12};
\node[ptitle] at (9.050,3.600) {(b) Stage 2, before the moves};
\node[pstat] at (9.050,0.000) {2 comm., $Q{=}0.239$, NMI 0.651};
\node[barrow, rotate=0, minimum height=9mm] at (12.283,1.800) {};
\node[note, anchor=south] at (12.283,2.080) {move 8};
\draw[goodc, line width=1.1pt, rounded corners=2pt] (12.933,0.000) rectangle (18.100,3.600);
\draw[ge] (17.850,2.297) -- (17.268,0.886);
\draw[ge] (17.850,2.297) -- (16.671,2.659);
\draw[ge] (17.850,2.297) -- (17.356,3.350);
\draw[ge] (17.850,2.297) -- (16.988,1.780);
\draw[ge] (15.652,2.836) -- (15.937,1.824);
\draw[ge] (15.652,2.836) -- (16.671,2.659);
\draw[ge] (15.652,2.836) -- (17.356,3.350);
\draw[ge] (15.652,2.836) -- (16.988,1.780);
\draw[ge] (15.652,2.836) -- (14.931,2.076);
\draw[ge] (15.937,1.824) -- (17.268,0.886);
\draw[ge] (15.937,1.824) -- (16.671,2.659);
\draw[ge] (15.937,1.824) -- (16.988,1.780);
\draw[ge] (15.937,1.824) -- (16.218,0.894);
\draw[ge] (17.268,0.886) -- (16.988,1.780);
\draw[ge] (16.671,2.659) -- (17.356,3.350);
\draw[ge] (16.671,2.659) -- (16.988,1.780);
\draw[ge] (17.356,3.350) -- (16.988,1.780);
\draw[ge] (16.988,1.780) -- (16.218,0.894);
\draw[ge] (16.218,0.894) -- (14.982,1.109);
\draw[ge] (14.982,1.109) -- (14.931,2.076);
\draw[ge] (14.982,1.109) -- (13.870,1.503);
\draw[ge] (14.931,2.076) -- (14.336,2.888);
\node[ring] at (16.218,0.894) {};
\node[gn, fill=pc0, text=white] (P3n1) at (17.850,2.297) {1};
\node[gn, fill=pc0, text=white] (P3n2) at (15.652,2.836) {2};
\node[gn, fill=pc0, text=white] (P3n3) at (15.937,1.824) {3};
\node[gn, fill=pc0, text=white] (P3n4) at (17.268,0.886) {4};
\node[gn, fill=pc0, text=white] (P3n5) at (16.671,2.659) {5};
\node[gn, fill=pc0, text=white] (P3n6) at (17.356,3.350) {6};
\node[gn, fill=pc0, text=white] (P3n7) at (16.988,1.780) {7};
\node[gn, fill=pc0, text=white] (P3n8) at (16.218,0.894) {8};
\node[gn, fill=pc1, text=white] (P3n9) at (14.982,1.109) {9};
\node[gn, fill=pc1, text=white] (P3n10) at (14.931,2.076) {10};
\node[gn, fill=pc1, text=white] (P3n11) at (14.336,2.888) {11};
\node[gn, fill=pc1, text=white] (P3n12) at (13.870,1.503) {12};
\node[ptitle] at (15.517,3.600) {(c) Synwalk move of node 8};
\node[pstat] at (15.517,0.000) {2 comm., $Q{=}0.207$, NMI 1.000};
\end{tikzpicture}
\caption{A Synwalk boundary move on a planted graph (nodes 1--8 and 9--12).}
\label{fig:m6}
\end{figure*}
